\section{Discussion}

\subsection{Consensus Is an Instrumental Achievement}

The central methodological parallel of this study is between annotation guidelines and Natural Language Inference (NLI) hypotheses. Detailed guidelines do not merely help human annotators express judgments that were already complete; they delimit which assumptions may be used, which ambiguities may be ignored, and which distinctions count as relevant. NLI hypotheses perform an analogous operation. A hypothesis conditions the model's activation pattern and allocation of attention, making some semantic relations salient while suppressing others. It may reduce uncontrolled variation, but this does not make the resulting judgment more neutral. In both cases, the instrument partly generates the regularity it records.

Consensus therefore cannot be treated as a synonym for semantic clarity. Agreement among annotators can be induced by a sufficiently restrictive codebook, just as agreement among NLI outputs can be induced by hypotheses that share presuppositions or lexical structure. The agreement may be useful, reproducible, and operationally necessary while still being an achievement of the probing conditions rather than an observer-free property of the text. The relevant question is not simply whether judgments converge, but under which constraints they converge, how long that convergence survives perturbation, and what must be excluded to maintain it. Consensus becomes an object to explain rather than a gold label at which explanation stops.

This reframing does not imply that every interpretation is equally warranted. It distinguishes the stability of an inference from the conditions that produced it. A stable judgment across genuinely different probes is stronger evidence than a stable judgment across paraphrases of one premise; a disagreement that persists across models, prompt forms, and corpora is more consequential than one caused by a single lexical accident. The aim is thus neither compulsory agreement nor unrestricted relativism, but an account of the transformations under which a semantic relation remains coherent.

\subsection{From Disagreement to Latent-Space Tomography}

Observer-based NLI offers a distinctive advantage because interpretive frames can be varied under controlled conditions while their outputs remain comparable within a shared representational construction. Rather than assigning one source label and then aggregating from it, the Semantic Interferometer admits each contrastive observer probe before downstream adjudication. This produces a form of latent-space tomography: the text is examined through a set of deliberately different semantic slices, and the repeated structure of their deformations is used to infer the shape that no single slice can reveal.

On this account, disagreement is not a residual error term. It can be described through coherence across perturbations, stability under resampling, directional asymmetry, local density, traversal cost, and perspectival shear. These quantities ask different questions. Coherence asks whether an observer-conditioned pattern repeats; stability asks whether it survives reasonable changes in the instrument; asymmetry asks whether transport from one frame to another differs from the reverse path; density asks how much neighboring evidence supports a region; path cost asks what resistance must be overcome to connect positions; and shear asks how sharply local relations deform under a change of frame. None is reducible to a majority vote. Together they shift the unit of analysis from a label attached to a text to the structured relation between interpretations of that text.

The architecture's failed designs are theoretically important here. Pre-labeled signed axes, raw embedding concatenation, and sequential Generalized Procrustes alignment each attempted, in different ways, to make encoded text submit to a convenient global frame. Their failures indicate that learned semantic representations are not infinitely malleable. Anisotropy, neighborhood structure, kernel choice, and the topology of a particular model constrain which transformations can be made without destroying relevant relations. This resistance rules out the strongest form of sociological fluidity: perspective does not mean that a text can be moved anywhere at no cost. Observer dependence and structural constraint coexist. The space may lack a privileged origin while retaining ridges, voids, barriers, and comparatively stable routes.

\subsection{Measurement, Worldmaking, and Theoretical Limits}

Goodman's account of worldmaking clarifies why this distinction matters. The stars can exist independently of a classification, while a constellation exists only through a practice that selects stars, draws relations among them, and gives the resulting figure a name. The constellation is not therefore useless or wholly arbitrary. It may guide navigation and organize shared experience. Yet it is a constituted object: the lines are supplied by a scheme of description rather than discovered as luminous filaments in the sky. Semantic measurement operates similarly. A left--right axis, an entailment hypothesis, a kernel, or a neighborhood rule does not simply expose an already bounded object called ``bias.'' It participates in deciding which differences belong to that object and which disappear outside it.

This is the deeper force of Goodman's 1978 argument for the present study. Measurement does not merely influence a final score; it helps constitute the measurable semantic object. Choosing observers creates a field of possible contrasts. Choosing a shared projection determines which relations can be compared. Choosing a density estimator decides what counts as support, and choosing a path rule decides what counts as proximity or resistance. The resulting geometry is neither a transparent copy of political language nor a free invention. It is a constrained construction produced by the interaction of texts, model parameters, probes, and analytic decisions.

Kuhn's 1962 account of theory-laden observation supplies the corresponding limit. There is no empirical output here without an instrument and no instrument without a framework of relevance. In this system, that framework includes at least the pretrained model, corpus composition, prompt or hypothesis bank, kernel family, random-feature basis, random seed, and downstream metric. Changing any of these can change the observable without changing the underlying article. Accordingly, reproducibility should not be understood as the recovery of one context-free map. It should mean disclosing provenance and testing which relations remain stable across a declared family of models, corpora, prompts, kernels, and seeds. Kuhn does not make comparison impossible; he makes the conditions of comparison impossible to ignore.

\subsection{Identity, Transport, and the Burden of a Global View}

Homotopy Type Theory (HoTT) provides a disciplined language for expressing this problem. In HoTT, identity need not be treated as an unexamined assertion that two presentations are literally the same. Identity can be represented by paths, and claims that dependent structure is preserved require transport along those paths. Applied cautiously to semantic geometry, this means that an article represented under observer frame $o$ and the ``same'' article represented under frame $o'$ should not automatically be identified coordinate by coordinate. A comparison requires a map $T_{o\to o'}$, together with evidence about what that map preserves. Equivalence across frames is therefore a result of successful transport, not a presupposition licensed by a shared document identifier.

Topos and sheaf theory make a related point through locality. Perspectives may be modeled as local sections defined over partially overlapping contexts. Where contexts overlap, their restrictions can be compared for compatibility. If the required compatibility and gluing conditions hold, local sections may determine a global section. Crucially, this formalism does not declare that global sections are impossible. Some claims may glue cleanly across observers, and some relations may be sufficiently invariant to support a global description. It requires, however, that globality be demonstrated rather than smuggled in through averaging. A mean of incompatible local sections is not, merely by being a mean, a view from nowhere.

These mathematical frameworks do not by themselves prove that natural language is perspectival, nor do they empirically establish that every linguistic identity must vary by observer. Their contribution is to change the burden of proof. If observer-conditioned representations are treated as local charts, then a supposedly neutral center must be shown to be canonical and natural across admissible observer maps. For observer spaces $X_o$ with proposed centers $c_o$, neutrality would require more than numerical centrality. At minimum, one would expect a relation such as
\begin{equation}
T_{o\to o'}(c_o)=c_{o'}
\end{equation}
for admissible transports, compatibility with composition, and independence from arbitrary choices of prompt order, basis, aggregation rule, and reference observer. Even this equivariance would not by itself establish uniqueness or normative neutrality; many coherent families of points can be transported. A privileged center would need a canonical construction that survives the relevant changes of frame without quietly reinstalling one frame as the standard.

This study supplies no such privileged point. What it supplies is a machinery for asking whether relations commute with transport, whether neighborhoods survive alignment, and where paths accumulate stress. The conceptual consequence is decisive: the first mathematical object should be the observer map, not an assumed origin. Positions acquire meaning through the transformations that connect them.

\subsection{What the Diagnostics Establish}

The empirical results support this relational interpretation, but only within their tested scope. Stronger claims require replication across kernels, control families, seeds, corpora, and metric bases.

The observer-recentering ablations give complementary controlled evidence. In a small synthetic robustness grid, coordinate translation alone passed $0$ of $9$ cells. The original focus-weighted local recomputation passed $3$ of $9$, while a uniform-weighted Random Kitchen Sinks local recomputation passed $7$ of $9$. These are not universal measurements of human semantics. They show, under planted synthetic structure, that moving the origin of an existing picture is not equivalent to recomputing geometry, shear, and density in an observer-conditioned frame. They also show that the result is sensitive to weighting, kernel, and seed. Observer dependence is thus best supported here as a conditional mechanism and ablation finding, not as an unqualified metaphysical fact extracted from nine trials.

The diagnostic labeled Naturality should be interpreted with similar care. It measures the survival of local neighborhood relations under an admissible orthogonal Procrustes transport. The available manuscript record does not preserve a single auditable aggregate value, so no exact score is asserted here. More importantly, orthogonal transport is designed to preserve Euclidean distance, making the diagnostic principally a consistency check on relational structure rather than independent proof that language is observer-dependent. Its relevance is conceptual: what survives the change of frame is a pattern of neighbors and relations, not evidence for one privileged origin. The natural object in this diagnostic is transportable structure. Naturality of relations does not entail naturality of a center.

Taken together, these findings support a restrained but consequential inference. There is structure in the observer-conditioned manifolds, some of that structure survives controlled transport, and meaningful recentering can require local recomputation. None of these results identifies an observer-free semantic basepoint. The evidence is better suited to an atlas of related charts than to a single coordinate plane declared neutral.

\subsection{The Geometric Fiction of the Neutral Middle}

This distinction sharpens the critique of systems such as Ground News. Dimensional reduction is not the problem. Any usable visualization must compress, and the Semantic Interferometer also produces rank-one and three-dimensional projections of much higher-dimensional objects. A projection becomes misleading when its construction is forgotten and one of its coordinates is reified as a property of the world. The error is not casting a shadow; it is calling the shadow the object and its midpoint the objective center.

Operationally, the Ground News middle is produced through a sequence of choices: source ratings supplied by external providers, placement on a left--right axis, aggregation and weighting rules, and an Overton Window that defines which positions count as politically legible. Each choice can be defensible for a product purpose. Together, however, they manufacture the geometry in which the middle appears. Treating that midpoint as an observer-free discovery converts a useful interface convention into a geometric fiction. This claim concerns the operation of the measurement system, not an assertion about the intentions of its designers.

From the HoTT perspective, no natural transport has been demonstrated that carries this middle unchanged across alternative source-rating schemes or observer frames. From the sheaf perspective, no compatibility argument has shown that local political judgments glue uniquely into that global point. From Goodman, the axis helps constitute the object it then reports; from Kuhn, the providers, categories, and political vocabulary are part of the instrument. A different provider set, issue taxonomy, weighting scheme, corpus, or historical window could move the center while leaving the articles untouched. The apparent neutrality of the middle is therefore an effect of suppressed provenance.

This does not mean that all placements are equally responsible or that factual reporting dissolves into opinion. A system can still test claims against evidence, expose source histories, detect fabrication, and distinguish well-supported interpretation from unsupported assertion. What it cannot legitimately do without further demonstration is elevate the barycenter of its chosen observers into a universal standpoint. The proper alternative is not to abolish comparison, but to display which observers produced the map, how the map changes under their replacement, and which relations remain invariant.

\subsection{Limits and a Revised Research Program}

The study's strongest philosophical conclusion must remain proportionate to these limits. The corpus is finite; parts of the validation are synthetic; the observer hypotheses are authored rather than inferred; the model inherits training-data regularities; the shared kernel is an approximation; and several results vary by kernel, seed, weighting, and control design. The latent topology belongs first to this model--corpus--probe system, not immediately to language as such. HoTT and topos theory organize the claims that would be required for cross-frame identity and globality, but they do not substitute for human validation, multilingual replication, model comparison, or adversarial testing.

Those limits point toward a more productive research program. Future systems should report observer and data provenance, uncertainty over maps, sensitivity to prompt and kernel choices, and the transformations under which conclusions persist. Evaluation should compare path structure, neighborhood stability, transport cost, and asymmetry across frames instead of rewarding only convergence toward a single label. In contested and interpretive NLP, the target should be an accountable geometry of differences, not the disappearance of difference into a centroid.

\sectionsection{Conclusion}

The ultimate implication of this study is the mathematical death of the objective center. This phrase does not announce the death of facts, evidence, semantic structure, or accountability. It names the failure of a more specific presumption: that every field of discourse contains a privileged, observer-free basepoint from which all other positions can be measured as deviations. Once semantic representations are understood as observer-conditioned charts connected by transports, the center ceases to be a discovery guaranteed by geometry. It becomes one constructed point among others unless its canonicity, naturality, and cross-frame stability are demonstrated.

The resulting picture is neither objectivist nor relativist in the ordinary sense. There are real constraints: learned spaces resist arbitrary deformation; some neighborhoods survive transport; some paths are costly; some claims glue across contexts; and evidence can rule interpretations in or out. Yet none of this requires a neutral origin. Structure resides in relations, invariants, obstructions, and paths. The mathematical death of the center therefore leaves more empirical work, not less. It demands that a system show how perspectives connect, where they fail to connect, and which features persist when the observer changes.

For systems like Ground News, the practical lesson is that a midpoint assembled from source raters, a left--right axis, aggregation weights, and a historically contingent Overton Window should be presented as the output of those choices. It may remain a useful summary, but it cannot stand in for an observer-free political middle. The geometry must expose its makers.

For NLP, the positive program is equally direct. Contested and interpretive tasks should stop trying to align bias away toward a supposed neutral centroid. They should instead measure the structured distance between perspectives, the transport required to move among them, the shear produced by that movement, and the path-dependent costs and asymmetries that remain. This recommendation does not reject factual or safety alignment in general. Models still require constraints against fabrication, harm, and unsupported claims. It rejects the narrower practice of confusing safety, truthfulness, or methodological discipline with the erasure of perspectival difference.

The fundamental unit of discourse is therefore not the isolated statement at a fixed coordinate, nor the majority label attached to it, but the provenance-bearing relation between situated interpretations. A mature semantic system would report that relation with uncertainty: which frame was used, which transport connected it to another, what changed along the path, what remained invariant, and where gluing failed. The replacement for the objective center is not an empty space. It is an atlas---a structured, testable account of the distances between perspectives. In that atlas, disagreement is no longer noise around truth's presumed location. It is part of the geometry through which truth claims become comparable at all.
